\section{Conclusions}
\label{sec:conclusions}

\status{Depends on the sections above. What can be written now is listed; the
rest is marked.}

\missing{conclusions --- to be written once the matrix is closed}

\subsection*{What is established}

\begin{itemize}
\item Both methods reproduce the ellipsoidal regime of the Newtonian reference
case and the shape of its published rise-velocity curve, and they \emph{converge
toward each other} under refinement --- from 1.94\% to 0.60\% in centre of mass
and from 4.3\% to 1.1\% in terminal rise velocity --- approaching from opposite
sides.
\item At the reference resolution \textbf{VOF} conserves mass about five times
better ($-0.170\%$ against $-0.851\%$), and front-tracking's drift \emph{grows}
with refinement. The coarse-mesh figures reverse this ranking and are partly
cancellation; they should not be quoted.
\item The two methods cost the same to within 5\% per cell-step, with the sign
of the difference depending on resolution.
\item The verification apparatus separates exact from convergent oracles, and
the exact ones --- surface closure, genus, the plane/slab/tetrahedron fraction
--- pass at machine precision.
\item The viscoelastic matrix reproduces the reference results. All twelve
ordering claims of the reference thesis hold, the quoted values agree to within
7\% (most within 3\%), the trend against a Newtonian reference in the same
geometry holds in all four cases, and the specific velocity signature the
reference describes --- a sharp drop after the maximum, then recovery ---
appears in both methods.
\item \textbf{VOF sits above the published values on all nine comparisons and
front-tracking at or below on seven of nine}, with mean absolute deviations of
$3.6\%$ and $1.2\%$. The asymmetry is the opposite of what the provenance
suggests, since the reference was computed with VOF.
\item Oldroyd-B did \emph{not} run away on the time scale of the rise: all
twelve runs completed 2001 steps with zero non-finite values.
\item \textbf{The cusp pair is closed at the reference resolution, and the two
methods agree on the integrated shape of the cusp to $0.5\%$} --- asymmetry
$1.2846$ in front-tracking against $1.2907$ in VOF, tracking each other to
$0.2$--$2.4\%$ from $t=1.6$ onward. This is the discriminating case, with the
two interface representations at their most different.
\item But their agreement in the \emph{local} curvature at $t=4$ is not
convergence: both give $R \approx 1.07\,h$, one cell. Two methods pinned at the
same mesh floor agree for a reason unrelated to the physics. What the column
does show is the \emph{approach}: VOF reaches its floor by $t=2.4$ and sits
there, front-tracking descends gradually and arrives only at $t=4$ --- the
signature of a per-cell planar reconstruction saturating at the cell while an
explicit surface keeps resolving below it.
\item At $h=1/25$ neither method had resolved the cusp, for two unrelated
reasons: VOF reports a trailing edge as round as the bubble ($R/R_{\mathrm{eq}}
= 0.71$), and the unfiltered front-tracking reports a spuriously sharp one
($0.047$) because its lowest vertex is a zig-zag spike. The filtered $h=1/50$
run gives $0.086$ --- \emph{less} sharp than the unfiltered coarse run, which is
the direct evidence that the coarse value was undulation.
\item The drop is $5.9\%$ at $h=1/25$ and $6.6\%$ at $h=1/50$ --- a bound below
the measurement's own noise rather than a resolved dependence.
\end{itemize}

\subsection*{What was withdrawn}

One result in this paper was retracted by its own measurement, and it is recorded
here rather than quietly dropped. An earlier version reported the velocity drop
as \emph{systematically larger in VOF} ($8.6$--$13.5\%$) than in front-tracking
($4.1$--$6.6\%$), and offered it as a difference between the two interface
representations. That number came from taking the global maximum of $w(t)$ on a
curve that is nearly flat after the overshoot --- so it was set by whichever
single frame sat highest, and in three of seven VOF runs it selected a frame at
$t \approx 3$--$4$, long after the physical overshoot. Locating the overshoot
physically and smoothing, front-tracking gives $5.5$--$8.1\%$ and VOF
$5.1$--$7.0\%$: the methods agree. What actually differs is the frame-to-frame
scatter, $3.0\%$ of the peak against $3.8$--$4.9\%$, and the raw definition
converted that scatter into an apparent physical asymmetry
(Section~\ref{sec:dropdef}).

The scatter is itself evidence on an open question below: a noisier extraction is
what the marching-cubes-bias reading predicts.

\subsection*{What is not}

\begin{itemize}
\item Adaptive refinement in production. The dynamic remeshing loop works and
is verified, but no case has been run on an adapted mesh to completion.
\item \emph{Resolved, and moved here from the open list.} Front-tracking's
mass drift appeared to grow with refinement because both production runs held
the surface subdivision fixed while the mesh was refined, so $ds/h$ doubled
between them. With $ds/h$ held constant the drift falls at an observed order of
$1.87$, and surgery reduces it by $42\%$ rather than causing it
(Section~\ref{sec:deriva-nsub}). What remains open is narrower: whether the
same holds at the production resolutions, which requires rerunning the coarse
case with the surface resolution matched to the grid.
\item Whether the $h=1/50$ agreement on the cusp survives separating resolution
from the filter. The coarse front-tracking run is unfiltered and the fine one is
filtered, so those two rows differ in two things at once. The cadence sweep of
Section~\ref{sec:tsur} produces a filtered coarse run as a by-product and
settles it.
\item Which of the two readings explains the sign asymmetry: a difference
between our volume-of-fluid and the reference's, or an upward bias in the
marching-cubes extraction from which the shape quantities are read.
\item Case~2, where topological change is expected.
\item Any statement about more than one process.
\end{itemize}
